\section{High-Frequency Probing}
\label{sec:method-probing}

Short local-return connections limited probe pickup and loop inductance.
Gate-source measurements require a source-referenced probe; switch-node
measurements require the full primary common-mode rating. Current-probe zero
and polarity were checked before acquisition. Voltage/current deskew is
required for switching-energy calculations. The CP2100A bandwidth of
\SI{800}{\kilo\hertz} is below the nominal \SI{1}{\mega\hertz} carrier,
limiting its use to input-current behavior rather than fast resonant-current
transitions.

A turn-on check needs switch-node voltage, the relevant gate or driver command,
and current direction through dead time to show whether the node reached the
required rail. The residual voltage at the intended turn-on instant is
\begin{equation}
 V_{\mathrm{res}}=
 \left|V_{\mathrm{SW}}(t_{\mathrm{on}})
 -V_{\mathrm{rail,target}}\right|.
 \label{eq:zvs-residual-method}
\end{equation}
The original prototype's burst-enable output PA10 marks the switching envelope,
not individual gate transitions. Enable and switch-node traces alone therefore
do not locate the ZVS measurement instant.

\begin{table}[htbp]
\centering
\caption{Probe and trigger settings needed for each test.}
\label{tab:probe-role-map}
\begin{tabularx}{\textwidth}{@{}p{0.27\textwidth}X@{}}
\toprule
Test & Probe and acquisition settings\\
\midrule
Startup/shutdown & Initial voltage, load, PA10 state, and trigger\\
Ripple/transient & Output node, probe coupling/bandwidth, sense point, and load event\\
No load/load removal & Disconnect method, output node, enable signal, and time span\\
Switching/ZVS & Gate/source reference, current direction, deskew, and bandwidth\\
Input power quality & Probe ranges, current zero/polarity, integration interval, and harmonic range\\
\bottomrule
\end{tabularx}
\end{table}
